\section{Weak Coupling and Asymptotic Freedom}
\label{sec:weak_coupling}

In this section, we address the ultraviolet stability problem. We analyze the theory in the regime where the lattice spacing \(a \to 0\) (\(\beta \to \infty\)) while keeping physical quantities fixed. The central task is to control the renormalization group flow and ensure the spectral gap scales correctly.

\subsection{Uniqueness via Renormalization Group}
We employ a rigorous renormalization group transformation \(\mathcal{R}\) to map the large-\(\beta\) theory to a sequence of effective actions. The uniqueness of the vacuum state and the control of the continuum limit rely on the stability of this flow.

\begin{theorem}[Weak Coupling Uniqueness (Conditional)]
Assume the flow of the effective action remains within the strict domain of attraction of the Gaussian fixed point (Hypothesis $H_{RG}$). Then for \(\beta > \beta_0\), the infinite-volume Gibbs measure is unique.
\end{theorem}

\subsection{Renormalization Group Analysis}
We utilize the Polchinski Flow Equation formalism to control the effective action as high-momentum modes are integrated out.

\begin{theorem}[Stability of the Flow]
Under the defined flow equations, the renormalized trajectory satisfies the Callan-Symanzik equation with a negative beta function \(\beta(g) < 0\) for small \(g\). This ensures asymptotic freedom and provides the bounds necessary for the continuum limit construction.
\end{theorem}

\subsection{Quantitative Homogenization}
A key technical hurdle is proving that the effective lattice operators converge to their continuum counterparts. We use Quantitative Homogenization to establish uniform ellipticity estimates for the effective block-spin action.

\begin{lemma}
Let \(A_{eff}\) be the effective diffusion matrix derived from the block-spin transformation. There exist constants \(0 < c \le C\) such that \(c I \le A_{eff} \le C I\) uniformly in the scaling limit. This ensures the spectral gap of the transfer matrix does not collapse faster than the canonical scaling.
\end{lemma}

\subsection{Universality}
The final step in the weak coupling analysis is establishing Universality. Irrelevant operators (dimension > 4) generated by the lattice regulator are shown to decay as powers of the lattice spacing \(a\).
\begin{equation}
\langle \mathcal{O}_{phys} \rangle_{lat} = \langle \mathcal{O}_{phys} \rangle_{cont} + O(a^2 \log a)
\end{equation}
This ensures that the specific choice of lattice action (provided it falls within the universality class) does not affect the existence of the mass gap in the continuum.


